% !TEX TS-program = pdflatex
% !TEX encoding = IsoLatin

\documentclass[a4paper,11pt,twoside,english]{toptesi}
\usepackage[latin1]{inputenc}
\usepackage[T1]{fontenc}
\usepackage{lmodern}
\usepackage[a-1b]{pdfx}
\usepackage{hyperref}
\usepackage{cite}
\usepackage{listings}
\usepackage{color}
\usepackage{url}
\usepackage{multirow}
\usepackage{amsmath}
\usepackage{verbatim}

% Definition of a new language for listings
%\lstdefinelanguage{P4}
%{
%	% list of keywords
%	alsoletter=\#,
%	morekeywords={
%		header\_type,
%		fields,
%		header,
%		parser,
%		extract,
%		return,
%		default,
%		next,
%		\#define,
%		latest,
%		select,
%		apply,
%		control,
%		table,
%		reads,
%		actions,
%		modify\_field,
%		action,
%		valid,
%		nop,
%		exact,
%		lpm,
%		if,
%		hit,
%		and,
%		extern,
%		void,
%		bit
%	},
%	sensitive=true, % keywords are not case-sensitive
%	morecomment=[l]{//}, % l is for line comment
%	morecomment=[s]{/*}{*/}, % s is for start and end delimiter
%	morestring=[b]" % defines that strings are enclosed in double quotes
%}
%
%\definecolor{eclipseBlue}{RGB}{42,0.0,255}
%\definecolor{eclipseGreen}{RGB}{63,127,95}
%\definecolor{eclipsePurple}{RGB}{127,0,85}
%
%\lstset{
%	language={P4},
%	basicstyle=\fontsize{8}{9}\ttfamily, % Global Code Style
%	captionpos=b, % Position of the Caption (t for top, b for bottom)
%	extendedchars=true, % Allows 256 instead of 128 ASCII characters
%	tabsize=2, % number of spaces indented when discovering a tab 
%	columns=fixed, % make all characters equal width
%	keepspaces=true, % does not ignore spaces to fit width, convert tabs to spaces
%	showstringspaces=false, % lets spaces in strings appear as real spaces
%	breaklines=true, % wrap lines if they don't fit
%	frame=trbl, % draw a frame at the top, right, left and bottom of the listing
%	%frameround=tttt, % make the frame round at all four corners
%	framesep=4pt, % quarter circle size of the round corners
%	numbers=left, % show line numbers at the left
%	numberstyle=\tiny\ttfamily, % style of the line numbers
%	commentstyle=\color{eclipseGreen}, % style of comments
%	keywordstyle=\color{eclipsePurple}, % style of keywords
%	stringstyle=\color{eclipseBlue}, % style of strings
%}

\hypersetup{%
    pdfpagemode={UseOutlines},
    bookmarksopen,
    pdfstartview={FitH},
    colorlinks,
    linkcolor={blue},
    citecolor={blue},
    urlcolor={blue}
  }

\newtheorem{osservazione}{Osservazione}% Standard LaTeX
\newtheorem{definition}{Definition}


\begin{document}
\english

\ateneo{Politecnico di Torino}
\corsodilaurea{Communications and Computer Networks Engineering}

\titolo{The title of my MSc thesis.}

\candidato{Your Name \textsc{Your Surname}}
\relatore{Prof.\ Name1 \textsc{Surname1}}
\secondorelatore{Prof.\ Name2 \textsc{Surname2}}
\sedutadilaurea{\textsc{Academic~Year} 202*-202*}
\logosede{Cover/Logoblu}

\CorsoDiLaureaIn{Master Degree course in\space}
\TesiDiLaurea{Master Degree Thesis}
\InName{in}
\CandidateName{Candidate}
\AdvisorName{Supervisors}


\frontespizio % make cover page

\ringraziamenti

Write here you acknowledgments. 

\abstract
Write here your abstract.

\indici

\mainmatter

\chapter{Introduction}
In this chapter, you must provide an overview of your thesis, written such that a non-engineer person (one of your parents) should be able to read and understand at high level.
\begin{itemize}
\item what is the addressed scenario and the addressed problem?
\item why the problem you addressed is interesting and/or practical relevant?
\item what is you main contribution? 
\item which methodology you used? (theory, simulation, experiments, literature overview)?
\item how the thesis is organized? which topic is covered in each chapter?
 \end{itemize}
 
 In the following the titles of chapters and sections are just examples.
 
 \section{Hints for writing}
 
 \begin{itemize}
 \item Use short sentences.
 \item Never cut \& paste a sentence from any website or paper. It is a copyright infringement and Politecnico anti-plagiarism system will detect it.
 \end{itemize}

\section{Latex hints}

\subsection{Latex examples}

According to what has been discussed in Sec.~\ref{sec:rel}, the following relation holds:
\begin{equation}\label{eq:1}
E=\dfrac{m}{2c^2}
\end{equation}
Consider that~\eqref{eq:1} is contradicting~\cite{ein05}.

With the above example, you have learned how to include math expressions, how to refer to equations, to other sections and to cite a related work. 

\subsection{Latex hints}
\begin{itemize}
\item never use \verb|\\ | to go to a new line. Instead leave al least one line empty for a new line. Leave latex to use the proper formatting for the paragraph.
\item use \verb|~| for all cases you need a space that cannot be split across multiple lines. I.e., ``\verb|Sec.~\ref{xxx}|'' and not ``\verb|Sec.~\ref{xxx}|'', or  ``\verb|according to~\eqref{xxx}|'' and not ``\verb|according to \eqref{xxx}|''.
\end{itemize}

\section{Plagiarism}

To be approved, your thesis must pass a plagiarism check. As a reminder:
\begin{itemize}
\item it is forbidden to copy verbatim one or more sentences  from any other external source (web pages, papers, etc.), unless you explicitly  grant the source and put the sentence/sentences in ``...''. Your can instead rephrase the whole sentence/sentences.
\item it is forbidden to include an image (diagram or figure) taken from web pages or papers. It is better to redraw. If you really need to include  an image from an external source, you must add to the caption ``(Reproduced from~\cite{example})'', citing properly the source.
\end{itemize}

 
\chapter{System model}

This chapter should provide all the required information for an engineer who does not know anything about your thesis topic and would like to get the required background to understand your work and contribution
 
\section{Model description}
...  

\section{Related work}\label{sec:rel}
After having described the system model, you should discuss about the related papers. Each work should be described specifying the scenario, the provided main contributions and you should highlight the similarities and differences with respect to your work in this thesis.

\chapter{My main contribution}\label{sec:cont}
This chapter should be devoted to describe your main contribution, e.g., the algorithms you devised/studied.

\chapter{Experimental/numerical evaluation}
Given the contribution in Chapter~\ref{sec:cont}, you must provide the adopted methodology and the  results

\section{Methodology}
Explain in all the details the adopted methodology, in order to make the results reproducible by anyone interested in addressing the same problem.

\section{Numerical results}

Here you will include all the main graphs. Each graph should appear with some explanation. The explanation must explain why the graph is interesting and what we learn from it, and not describing the numerical values you have on it (the reader is not blind and see by himself the numbers). Consider the example of Fig.~\ref{fig:bw}.

\begin{figure}[tb!]
\begin{center}
\includegraphics[width=10cm]{Figures/sino.pdf}
\caption{This is the required caption of the graph.}\label{fig:bw}
\end{center}
\end{figure}

Some basic hints:
\begin{itemize}
\item both x and y axes should be properly labeled. Units of measurements must be {\em always} reported!
\item if the curve are experimental or by simulation, remember to report always all the points in addition to the line connecting them
\end{itemize}




\chapter{Conclusion}

Here you must summarize completely your work. You do not need anymore to motivate your topic, but just report in short what you have done and the main results.

If you wish, you can also add the possible direction for future works, extending your thesis work.

\bibliographystyle{plain}
\bibliography{biblio}

\end{document}
